\documentclass{situationist-preprint}

\title{The Pleasure of Tragedy}
\msid{AES-06-a}
\msfamily{AES-06}
\msdiscipline{Aesthetics}
\mstype{Principal manuscript}
\mskeywords{tragedy; paradox of tragedy; negative emotions; catharsis; understanding; Hume; Aristotle}

\begin{document}
\maketitle

\begin{abstract}
Why do we enjoy tragedy? Tragedies elicit pity, fear and grief, which are painful, and people seek them out. We argue that tragedy offers a controlled engagement with loss, and that its pleasure is a form of understanding. The engagement is controlled because the emotions it elicits are directed at fictions and bounded by the frame of the work; the pleasure lies not in the painful emotions themselves but in the understanding of loss that they make possible and that could not be achieved without them. We state the paradox as an inconsistent triad, show that our account rejects the third member only in its unrestricted form, and compare it with the conversion theory of Hume, the meta-response theory of Feagin and the view that negative emotions are themselves enjoyed. The account explains two features the rivals do not: that the pleasure of tragedy is reflective rather than immediate, and that it is diminished when a tragedy is felt to be gratuitous.
\end{abstract}

\section{The paradox}

Aristotle observed that we take pleasure in contemplating accurate images of things we find painful to see in reality, such as the forms of the meanest animals and of corpses, and explained it by the pleasure of learning (\emph{Poetics} 1448b). Tragedy is the extreme case. Its proper effect, in his account, is to arouse pity and fear (\emph{Poetics} 1449b), and audiences return to it willingly. Hume opened his essay on the subject by noting the \enquote{unaccountable pleasure} that spectators of a well-written tragedy receive from sorrow, terror and anxiety \citep{Hume1757}. The puzzle can be stated as three claims that seem individually plausible and jointly inconsistent.

\begin{enumerate}[label=(T\arabic*),nosep]
\item People take pleasure in tragedies.
\item Tragedies elicit painful emotions, and elicit them essentially: a tragedy that did not would fail as a tragedy.
\item People do not take pleasure in what elicits painful emotions.
\end{enumerate}

Solutions deny one of the three. Some deny (T2), holding that the emotions elicited by fiction are not really painful, or not really emotions. Most deny (T3) in some form. They differ on why it fails.

\section{Positions}

\emph{Conversion.} \citet{Hume1757} argued that the eloquence and art of the representation convert the painful passions into pleasure: the predominant emotion, delight in the beauty of the work, absorbs the subordinate ones and draws their force into itself. The difficulty, pressed since Hume, is that the mechanism is obscure, and that the pity and fear of the audience do not seem to be converted into anything; they remain pity and fear, and the audience leaves moved, not merely delighted.

\emph{Compensation.} On compensatory views, the painful emotions remain painful, and the pleasure comes from something else that outweighs them: the beauty of the language, the satisfaction of a well-made plot, the exercise of the imagination. Such views explain why we accept the pain but not why tragedy should require it. If the pleasure lay in the language, a beautifully written comedy would serve as well.

\emph{Meta-response.} \citet{Feagin1983} distinguished a direct response to the events of a tragedy, which is painful, from a meta-response to the direct response, which is pleasant: we find that we are the kind of people who are moved by the suffering of others, and approve of it. The meta-response explains why the pain is necessary, since there must be a direct response to approve of. It makes the pleasure of tragedy a form of self-approval, which seems too self-regarding to capture what audiences value.

\emph{Enjoyment of negative emotion.} \citet{Morreall1985} and others hold that negative emotions can themselves be enjoyed when we are in control of the situation that produces them. This denies (T3) directly. It is supported by the observation that people seek frightening rides and sad music, but it does not explain why tragedy is valued more highly than a frightening ride, or why its pleasure seems to be one of insight.

\emph{Resignation and affirmation.} The nineteenth century gave the paradox a metaphysical turn. \citet{Schopenhauer1819} held that tragedy shows the vanity of willing: the spectacle of the will turned against itself, in the suffering of the hero, leads the spectator to resignation, and the pleasure of tragedy is the relief of that detachment. \citet{Nietzsche1872} reversed the valuation. Tragedy does not counsel resignation but affirms life in the face of suffering, by letting the spectator experience, beneath the destruction of the individual, the indestructible vitality of the whole. Both accounts make the pleasure of tragedy depend on a metaphysics that few now accept. Both nevertheless identify something the earlier accounts missed: the pleasure of tragedy is bound up with a change in how the spectator sees the place of suffering in human life. The present account keeps this insight without the metaphysics.

\emph{Distancing and embracing.} Recent empirical work proposes that art reception involves both psychological distancing, which makes negative emotions tolerable, and embracing, through which they are integrated into a rewarding experience \citep{MenninghausEtAl2017}. The model is compatible with several philosophical accounts. It supports the claim that control is a precondition of the pleasure, without saying what the pleasure is.

\section{Framework}

\begin{definition}\label{d:tragedy}
A \emph{tragedy}, in the sense relevant here, is a representation of a serious loss suffered by persons with whom the audience is invited to sympathize, constructed so that the loss follows intelligibly from the characters and their situation.
\end{definition}

The last clause is Aristotle's requirement that the events follow by necessity or probability (\emph{Poetics} 1451a). It distinguishes tragedy from the mere depiction of suffering.

\begin{principle}[Control]\label{p:control}
The emotions elicited by a tragedy are bounded: they are directed at fictional events, they lack the practical consequences of emotions directed at real events, and their duration is fixed by the work.
\end{principle}

AES-05 argues that the emotions elicited by fiction are quasi-emotions, which share the physiological profile of emotions without the corresponding beliefs. \ref{p:control} is the consequence of that account for tragedy: the audience's pity is felt but does not issue in attempts to help, and its fear is felt but does not issue in flight.

\begin{principle}[Understanding through emotion]\label{p:understanding}
Some understanding of loss is available only through an emotional response to it: to understand what a loss means for a person is in part to feel its weight.
\end{principle}

\ref{p:understanding} is a version of the view that emotions are a mode of cognition, defended by \citet{Nussbaum1986} for tragedy in particular. It does not claim that emotions are infallible, only that some evaluative understanding cannot be achieved by a subject who feels nothing.

\section{The argument}

\begin{proposition}\label{p:necessary}
The painful emotions elicited by a tragedy are necessary for the understanding it offers.
\end{proposition}

\begin{proof}
By Definition~\ref{d:tragedy}, a tragedy represents a serious loss as following intelligibly from its circumstances, so what it offers to understand is a loss and its weight for those who suffer it. By \ref{p:understanding}, that understanding requires an emotional response to the loss, and the response appropriate to a serious loss suffered by persons one sympathizes with is pity, fear or grief.
\end{proof}

\begin{theorem}\label{t:main}
The pleasure of tragedy is the pleasure of understanding a loss, achieved through painful emotions whose bounds make them tolerable.
\end{theorem}

\begin{proof}
Understanding is a source of pleasure, as Aristotle observed of learning from imitation. By Proposition~\ref{p:necessary}, the understanding tragedy offers requires painful emotions. By \ref{p:control}, these emotions are bounded, so the cost of the understanding is limited in a way that the cost of understanding a real loss is not. A tragedy therefore provides an understanding of loss at a price an audience can choose to pay.
\end{proof}

The theorem resolves the paradox by restricting (T3). People do not take pleasure in what elicits painful emotions as such, but they can take pleasure in an understanding that is available only through painful emotions, when the emotions are bounded. This is the sense in which tragedy offers a controlled engagement with loss. It also gives the old idea of catharsis a defensible reading. On the interpretation of \citet{Golden1962}, catharsis is clarification, an intellectual rather than a purgative process, and this is close to the present account. The companion manuscript AES-06-b examines the interpretive question.

Two features of the experience of tragedy support the account over its rivals. First, the pleasure of tragedy is reflective rather than immediate: it is felt most strongly at the end of the work and afterwards, when the shape of the loss has become clear, and not at the moments of greatest pain. This is predicted if the pleasure is one of understanding, and not if it lies in the negative emotions themselves or in the beauty of the language, which is present throughout. Second, a tragedy whose catastrophe is felt to be gratuitous, not following from the characters and their situation, is experienced as merely distressing. This is predicted by Definition~\ref{d:tragedy} and Proposition~\ref{p:necessary}: where the loss is not intelligible, there is nothing to understand, and the painful emotions are pure cost.

\section{Further cases}

\emph{Oedipus.} \emph{Oedipus Tyrannus} is the case on which Aristotle built his account, and it shows the role of intelligibility. Oedipus's catastrophe follows from his own investigation: the qualities that made him a ruler, his determination to know and to act on what he knows, are the qualities that destroy him. The audience's pity and fear are directed at a loss whose shape is fully visible by the end, and the pleasure of the play is concentrated in the recognition scene, where the shape becomes clear. This is what Theorem~\ref{t:main} predicts: the pleasure peaks at the moment of greatest understanding, which is also a moment of great pain.

\emph{King Lear.} Lear is often taken as a counterexample to any cognitive account, since the death of Cordelia is notoriously gratuitous; Johnson found it unbearable, and for a century and a half the play was performed in Tate's version, in which she survives. The case confirms the account rather than refuting it. The ending was felt as unbearable precisely because it seemed unintelligible, and the critical tradition that restored it did so by finding it intelligible: as the completion of Lear's discovery of what he had lost, which is incomplete until the loss is total. Where readers find the ending intelligible, it becomes the source of the play's distinctive power; where they do not, it remains a wound. The account predicts this division of response.

\emph{Sad music.} Instrumental music that expresses sadness poses a structurally similar puzzle, which \citet{Levinson1982} discussed under the heading of the negative emotions in music. AES-07 argues that music expresses emotions by resembling their characteristic behavior. Whether sad music offers understanding of loss is doubtful, since it has no represented loss to understand. The present account therefore does not extend to it without modification. Its pleasure may be closer to the enjoyment of negative emotion as such, or to the pleasure of recognizing expressive form, and the difference between the experience of sad music and that of tragedy is evidence that the two pleasures differ.

\section{Objections and replies}

\begin{objection}
Many people enjoy tragedies they have seen before, whose outcome they know. There is nothing new to understand.
\end{objection}

\begin{reply}
Understanding the weight of a loss is not the same as knowing that it occurs. One can know the outcome of \emph{King Lear} and still come to understand, through a new performance, what Lear's loss consists in. In the same way, a person can know that a friend has died and still be coming to understand, months later, what the death means. The repetition of tragedy is like the rereading of a philosophical argument one already knows: the conclusion is familiar, and the understanding deepens.
\end{reply}

\begin{objection}
Some people enjoy sad stories that offer no understanding at all: melodramas whose losses are contrived and whose emotions are cheap.
\end{objection}

\begin{reply}
The objection is correct, and the account explains the difference in value between tragedy and melodrama. The enjoyment of melodrama is closer to the enjoyment of negative emotion as such, which Morreall describes. It exists, and it is a lesser pleasure, because it lacks the understanding that tragedy adds. Our thesis concerns tragedy, and the contrast with melodrama is evidence for it: what tragedy has and melodrama lacks is intelligibility of the loss, and that is what the account identifies as the source of the distinctive pleasure.
\end{reply}

\begin{objection}
Understanding is not always pleasant. Understanding that one has been betrayed is painful.
\end{objection}

\begin{reply}
When the object of understanding is a loss of one's own, the pain of the loss swamps the pleasure of understanding it. \ref{p:control} is what makes the difference in tragedy: the loss is not one's own, and the emotions it elicits are bounded, so the pleasure of understanding is not swamped. This is also why the pleasure of tragedy can fail for a spectator whose own situation is too close to the one represented: the bounds of \ref{p:control} no longer hold for her.
\end{reply}

\section{What the result does not establish}

The account does not show that every pleasure taken in tragedy is a pleasure of understanding; audiences also enjoy spectacle, language and performance, and these pleasures are real. It identifies the pleasure that is distinctive of tragedy and that explains why tragedy requires painful emotions. Nor does it settle the question, raised by \citet{Smuts2009}, whether there are artworks valued for negative affect that offers no understanding; horror may be such a case, though \citet{Carroll1990} argues that its pleasure is also cognitive, lying in the satisfaction of curiosity about the impossible.

\section{Conclusion}

We enjoy tragedy because it lets us understand loss, and the understanding is available only through pity, fear and grief. Tragedy bounds these emotions, so that the understanding can be had at a cost the audience is willing to bear. The paradox arises from an unrestricted principle that people do not seek what pains them; restricted to unbounded pain, the principle is true, and tragedy does not violate it.

\begin{thebibliography}{99}\small
\bibitem[Carroll(1990)]{Carroll1990} Carroll, N. (1990). \emph{The Philosophy of Horror, or Paradoxes of the Heart}. Routledge.
\bibitem[Feagin(1983)]{Feagin1983} Feagin, S.~L. (1983). The pleasures of tragedy. \emph{American Philosophical Quarterly}, 20(1), 95--104.
\bibitem[Golden(1962)]{Golden1962} Golden, L. (1962). Catharsis. \emph{Transactions and Proceedings of the American Philological Association}, 93, 51--60.
\bibitem[Hume(1757)]{Hume1757} Hume, D. (1757). Of tragedy. In \emph{Four Dissertations}. Millar.
\bibitem[Levinson(1982)]{Levinson1982} Levinson, J. (1982). Music and negative emotion. \emph{Pacific Philosophical Quarterly}, 63(4), 327--346.
\bibitem[Menninghaus et~al.(2017)]{MenninghausEtAl2017} Menninghaus, W., Wagner, V., Hanich, J., Wassiliwizky, E., Jacobsen, T. and Koelsch, S. (2017). The distancing--embracing model of the enjoyment of negative emotions in art reception. \emph{Behavioral and Brain Sciences}, 40, e347.
\bibitem[Morreall(1985)]{Morreall1985} Morreall, J. (1985). Enjoying negative emotions in fictions. \emph{Philosophy and Literature}, 9(1), 95--103.
\bibitem[Nietzsche(1872)]{Nietzsche1872} Nietzsche, F. (1872). \emph{Die Geburt der Tragödie aus dem Geiste der Musik}. Fritzsch.
\bibitem[Nussbaum(1986)]{Nussbaum1986} Nussbaum, M.~C. (1986). \emph{The Fragility of Goodness: Luck and Ethics in Greek Tragedy and Philosophy}. Cambridge University Press.
\bibitem[Schopenhauer(1819)]{Schopenhauer1819} Schopenhauer, A. (1819). \emph{Die Welt als Wille und Vorstellung}. Brockhaus.
\bibitem[Smuts(2009)]{Smuts2009} Smuts, A. (2009). Art and negative affect. \emph{Philosophy Compass}, 4(1), 39--55.
\end{thebibliography}

\end{document}
